\subsection{在李代数中约化的子代数}

定义 5. 设 g 是李代数，h 是 g 的李子代数。如果表示 \(x \mapsto ad_{g} x\) 给出的 h 在 g 上的表示是半单的，则称 h 在 g 中是约化的。

该表示包含 \(\mathfrak{h}\) 的伴随表示作为子表示。因此，若 \(\mathfrak{h}\) 在 \(\mathfrak{g},\mathfrak{h}\) 中是约化的，则它是约化的。另一方面，说一个李代数在自身中约化，等价于说它是约化的。

命题 7. 设 g 是李代数，h 是在 g 中约化的子代数，ρ 是 g 在向量空间 V 上的一个表示，W 是 V 中所有作为有限维单 h-模的子空间之和。则 W 在 ρ(g) 下稳定。

令 \(W_0\) 为 V 的有限维单 \(\mathfrak{h}\)-子模。我们要证明关系 \(\rho(x)(W_0) \subset W\) 对所有 \(x \in \mathfrak{g}\) 成立。令 M 表示向量空间 \(\mathfrak{g}\)，将其视为 \(\mathfrak{h}\)-模，所依据的是 \(x \mapsto \operatorname{ad}_{\mathfrak{g}} x\) 给出的 \(\mathfrak{h}\) 在 \(\mathfrak{g}\) 上的表示。于是 \(M \otimes_{K} W_0\) 是半单 \(\mathfrak{h}\)-模（定理 4 的推论 1）。令 \(\theta\) 为从 \(\mathbf{M} \otimes_{\mathbb{K}} \mathbf{W}_0\) 到 \(\mathbf{V}\) 的 K-线性映射，由 \(\theta(x \otimes w) = \rho(x)w\) 定义。它是 \(\mathfrak{h}\)-模同态，因为若 \(y \in \mathfrak{h}\) 则：

\[
\begin{array}{r l} \theta ([ y, x ] \otimes w + x \otimes \rho (y) w) & = \rho ([ y, x ]) w + \rho (x) \rho (y) w \\ & = \rho (y) \rho (x) w = \rho (y) \theta (x \otimes w). \end{array}
\]

因此 \(\theta (\mathbf{M}\otimes_{\mathbb{K}}\mathbf{W}_0)\) 是有限维半单 \(\mathfrak{h}\)-模。因此 \(\theta (\mathbf{M}\otimes_{\mathbb{K}}\mathbf{W}_0)\subset \mathbf{W}\) 包含于 \(\rho (x)(\mathbf{W}_0)\subset \mathbf{W}\)，也就是说对所有 \(x\in \mathfrak{g}\) 都有相应的包含关系。

推论 1. 设 g 是李代数，h 是在 g 中约化的子代数，\(\rho\) 是 g 的有限维半单表示。则 \(\rho\) 在 h 上的限制是半单的。

只需考虑 \(\rho\) 是单的情形。采用命题 4 的记号 V、W。令 \(W_{1}\) 为 V 中在 \(\rho(\mathfrak{h})\) 下稳定的非零子空间中极小的一个。则 \(W_{1} \subset W\)，从而 \(W \neq \{0\}\)，进而 \(W = V\)。

推论 2. 设 g 是李代数，h 是在 g 中约化的子代数，k 是 h 中在 h 内约化的子代数。则 k 在 g 中约化。

表示 \(x \mapsto \operatorname{ad}_{\mathfrak{g}} x\) 给出的 \(\mathfrak{h}\) 在 \(\mathfrak{g}\) 上的表示是半单的，因此它在 \(k\) 上的限制是半单的（推论 1）。

